%───────────────────────────────────────────────────────────────────────────────
% TCOLORBOX.TEX — Estilo base y entornos auxiliares para cajas
%───────────────────────────────────────────────────────────────────────────────
% Define basebox: estilo heredado por todos los entornos de environments/.
% Define boxfig y boxtab: versiones no-flotantes para usar DENTRO de cajas.
%
% PROBLEMA FUNDAMENTAL — flotantes dentro de tcolorbox:
%   tcolorbox construye internamente un \vbox. Cualquier entorno flotante
%   (\begin{figure}, \begin{table}, \begin{acadtable}) invoca \@float, que
%   requiere "outer par mode". Dentro de un \vbox ese modo no existe →
%   error "Not in outer par mode".
%
%   Entornos afectados dentro de cualquier caja basada en basebox:
%     \begin{figure}      → usar \begin{boxfig}
%     \begin{table}       → usar \begin{boxtab}
%     \begin{acadtable}   → usar \begin{boxtab}
%     \begin{lstlisting}[float=...] → quitar la opción float=
%     \begin{longtable}   → incompatible; usar tabularx sin float
%
%   La solución usa \captionof{tipo}{texto} del paquete caption (ya cargado),
%   que genera el caption y actualiza el contador sin invocar el mecanismo
%   de flotante. \phantomsection corrige el anchor del PDF con hyperref.
%
% ANCHO DENTRO DE CAJAS:
%   Dentro de tcolorbox, \linewidth = ancho del área de contenido de la caja
%   (descontando left=6pt + right=6pt). Usar siempre \linewidth, no \textwidth,
%   para tabularx, \includegraphics[width=...], etc.
%───────────────────────────────────────────────────────────────────────────────


% ─── Estilo base: heredado por todos los entornos teoremáticos ───────────────
\tcbset{
  basebox/.style={
    enhanced,
    breakable,
    colback    = surface0,
    colframe   = overlay0!30,
    coltext    = maintext,
    arc        = 0.5pt,
    outer arc  = 0.5pt,
    boxrule    = 0.5pt,
    left       = 6pt,
    right      = 6pt,
    top        = 8pt,
    bottom     = 8pt,
    drop fuzzy shadow = background,
    fonttitle  = \bfseries,
    before title = {\vspace{2pt}},
    after title  = {\vspace{2pt}},
  }
}


% ─── boxfig: figura no-flotante para usar dentro de cajas ────────────────────
% Uso:
%   \begin{boxfig}{fig:label}{Caption de la figura}
%     \includegraphics[width=\linewidth]{ruta/imagen}
%   \end{boxfig}
%
% Nota: el label va DESPUÉS de \captionof para que hyperref ancle
% correctamente. \phantomsection asegura el anchor del PDF.
\NewDocumentEnvironment{boxfig}{ m m }{%
  \begin{center}
}{%
  \phantomsection
  \captionof{figure}{#2}%
  \label{#1}%
  \end{center}
}


% ─── boxtab: tabla no-flotante para usar dentro de cajas ─────────────────────
% Reemplaza \begin{acadtable} y \begin{table} dentro de tcolorbox.
%
% Uso:
%   \begin{boxtab}{tab:label}{Caption de la tabla}
%     \begin{tabularx}{\linewidth}{l Y Z}
%       ...
%     \end{tabularx}
%   \end{boxtab}
%
% Nota: usar \linewidth (no \textwidth) para tabularx dentro de la caja.
% La posición del caption sigue la convención de tablas: arriba del contenido.
\NewDocumentEnvironment{boxtab}{ m m }{%
  \begin{center}
  \phantomsection
  \captionof{table}{#2}%
  \label{#1}%
}{%
  \end{center}
}


% ─── Entornos anidados: cajas dentro de cajas ────────────────────────────────
% tcolorbox permite anidar cajas solo si la caja interna tiene breakable=false.
% Estilo para cajas secundarias dentro de un entorno teoremático:
\tcbset{
  innerbox/.style={
    basebox,
    breakable = false,        % obligatorio para anidar dentro de basebox
    boxrule   = 0.3pt,
    left      = 4pt,
    right     = 4pt,
    top       = 4pt,
    bottom    = 4pt,
    drop fuzzy shadow = none, % sin sombra en cajas internas
  }
}
